% GENERATED FILE. Edit canonical lessons, then run npm run generate:books.
% canonical-source-hash: fnv1a64:5ad8b1d1ed7c9c7d
% canonical-lessons: GE-C313-maerchen, GE-C313-legende, GE-C313-fabel, GE-C313-hexe, GE-C313-zauberer

\chapter{Es war einmal}
\label{ch:ge-es-war-einmal}
\hlchaptermodality{313}

\begin{chapteropening}
\textbf{By the end of this chapter:} \emph{I can talk about old tales: a fairy tale, a legend, a fable, a witch and a wizard.}

Complete the last lesson of chapter 313: I can talk about old tales: a fairy tale, a legend, a fable, a witch and a wizard.
\end{chapteropening}

\section[das Märchen]{das Märchen — the fairy tale}

\label{lesson:GE-C313-maerchen}



\begin{quote}
Before the new one: say the German for the traditional dress, then the German for the band.
\end{quote}

\subsection*{You'll want to know: das Märchen}
\textbf{das Märchen} — \textquotedblleft{}the fairy tale\textquotedblright{}.

\begin{quote}
Meine Oma erzählte uns jeden Abend ein Märchen.
\end{quote}

\begin{quote}
\emph{My grandma told us a fairy tale every evening.}
\end{quote}

Plural: \emph{die Märchen}. They begin \emph{Es war einmal} — once upon a time.

\begin{cousinweb}
\emph{Märchen} is the diminutive of the old word \emph{Mär} \textquotedblleft{}news, tale\textquotedblright{}, from Middle High German \emph{mære} (EtymDB), with the ending \emph{-chen}: a little tale.
\end{cousinweb}

\begin{grammarlens}[title={What you've built}]
One more word for telling, explaining and describing at B1.
\end{grammarlens}

\subsection*{Your turn}
\begin{itemize}
\raggedright
  \item \emph{Say it:} \emph{das Märchen}
  \item \emph{Say it:} \emph{das Märchen}, once more
  \item \emph{Say it:} a sentence of your own with \emph{das Märchen}
  \item \emph{Recall:} say the German for the traditional dress, then the German for the band, then say \emph{das Märchen} again
\end{itemize}

\begin{tcolorbox}[breakable,colback=teal!4,colframe=teal!35!black,title={Before you move on}]
What does das Märchen mean? (\textquotedblleft{}The fairy tale\textquotedblright{}.) Say it once more.
\end{tcolorbox}
\section[die Legende]{die Legende — the legend}

\label{lesson:GE-C313-legende}



\begin{quote}
Before the new one: say the German for the band, then the German for the fairy tale.
\end{quote}

\subsection*{You'll want to know: die Legende}
\textbf{die Legende} — \textquotedblleft{}the legend\textquotedblright{}.

\begin{quote}
Eine alte Legende erzählt von dieser Brücke.
\end{quote}

\begin{quote}
\emph{An old legend tells of this bridge.}
\end{quote}

Plural: \emph{die Legenden}.

\begin{cousinweb}
\emph{Legende} comes from Middle High German \emph{legende} \textquotedblleft{}the story of a saint\textquotedblright{} (EtymDB), from Medieval Latin \emph{legenda} \textquotedblleft{}things to be read\textquotedblright{}, built on \emph{legere} \textquotedblleft{}to read\textquotedblright{} (Lewis \& Short).
\end{cousinweb}

\begin{grammarlens}[title={What you've built}]
One more word for telling, explaining and describing at B1.
\end{grammarlens}

\subsection*{Your turn}
\begin{itemize}
\raggedright
  \item \emph{Say it:} \emph{die Legende}
  \item \emph{Say it:} \emph{die Legende}, once more
  \item \emph{Say it:} a sentence of your own with \emph{die Legende}
  \item \emph{Recall:} say the German for the band, then the German for the fairy tale, then say \emph{die Legende} again
\end{itemize}

\begin{tcolorbox}[breakable,colback=teal!4,colframe=teal!35!black,title={Before you move on}]
What does die Legende mean? (\textquotedblleft{}The legend\textquotedblright{}.) Say it once more.
\end{tcolorbox}
\section[die Fabel]{die Fabel — the fable}

\label{lesson:GE-C313-fabel}



\begin{quote}
Before the new one: say the German for the fairy tale, then the German for the legend.
\end{quote}

\subsection*{You'll want to know: die Fabel}
\textbf{die Fabel} — \textquotedblleft{}the fable\textquotedblright{}.

\begin{quote}
In der Fabel spricht der Fuchs mit der Maus.
\end{quote}

\begin{quote}
\emph{In the fable the fox talks to the mouse.}
\end{quote}

Plural: \emph{die Fabeln}. Animals that talk, and a lesson at the end.

\begin{cousinweb}
\emph{Fabel} comes from Middle High German \emph{fabel} (EtymDB), from Latin \emph{fābula} \textquotedblleft{}a story, a tale\textquotedblright{} (Lewis \& Short).
\end{cousinweb}

\begin{grammarlens}[title={What you've built}]
One more word for telling, explaining and describing at B1.
\end{grammarlens}

\subsection*{Your turn}
\begin{itemize}
\raggedright
  \item \emph{Say it:} \emph{die Fabel}
  \item \emph{Say it:} \emph{die Fabel}, once more
  \item \emph{Say it:} a sentence of your own with \emph{die Fabel}
  \item \emph{Recall:} say the German for the fairy tale, then the German for the legend, then say \emph{die Fabel} again
\end{itemize}

\begin{tcolorbox}[breakable,colback=teal!4,colframe=teal!35!black,title={Before you move on}]
What does die Fabel mean? (\textquotedblleft{}The fable\textquotedblright{}.) Say it once more.
\end{tcolorbox}
\section[die Hexe]{die Hexe — the witch}

\label{lesson:GE-C313-hexe}



\begin{quote}
Before the new one: say the German for the legend, then the German for the fable.
\end{quote}

\subsection*{You'll want to know: die Hexe}
\textbf{die Hexe} — \textquotedblleft{}the witch\textquotedblright{}.

\begin{quote}
Die Hexe wohnte in einem kleinen Haus im Wald.
\end{quote}

\begin{quote}
\emph{The witch lived in a small house in the forest.}
\end{quote}

Plural: \emph{die Hexen}.

\begin{cousinweb}
\emph{Hexe} comes from Middle High German \emph{hecse} (EtymDB), and EtymDB lists English \emph{hag} as a cousin, from Proto-Germanic \emph{hagatusjōn} \textquotedblleft{}witch\textquotedblright{}.
\end{cousinweb}

\begin{grammarlens}[title={What you've built}]
One more word for telling, explaining and describing at B1.
\end{grammarlens}

\subsection*{Your turn}
\begin{itemize}
\raggedright
  \item \emph{Say it:} \emph{die Hexe}
  \item \emph{Say it:} \emph{die Hexe}, once more
  \item \emph{Say it:} a sentence of your own with \emph{die Hexe}
  \item \emph{Recall:} say the German for the legend, then the German for the fable, then say \emph{die Hexe} again
\end{itemize}

\begin{tcolorbox}[breakable,colback=teal!4,colframe=teal!35!black,title={Before you move on}]
What does die Hexe mean? (\textquotedblleft{}The witch\textquotedblright{}.) Say it once more.
\end{tcolorbox}
\section[der Zauberer]{der Zauberer — the wizard, the magician}

\label{lesson:GE-C313-zauberer}



\begin{quote}
Before the new one: say the German for the fable, then the German for the witch.
\end{quote}

\subsection*{You'll want to know: der Zauberer}
\textbf{der Zauberer} — \textquotedblleft{}the wizard, the magician\textquotedblright{}.

\begin{quote}
Der Zauberer hatte einen langen weißen Bart.
\end{quote}

\begin{quote}
\emph{The wizard had a long white beard.}
\end{quote}

Plural: \emph{die Zauberer}. At a party, a magician.

\begin{cousinweb}
\emph{Zauberer} is built on \emph{Zauber} \textquotedblleft{}magic\textquotedblright{}, Old High German \emph{zoubar}, from Proto-Germanic \emph{taubrą} \textquotedblleft{}magic\textquotedblright{} (EtymDB).
\end{cousinweb}

\begin{grammarlens}[title={What you've built}]
One more word for telling, explaining and describing at B1.
\end{grammarlens}

\subsection*{Your turn}
\begin{itemize}
\raggedright
  \item \emph{Say it:} \emph{der Zauberer}
  \item \emph{Say it:} \emph{der Zauberer}, once more
  \item \emph{Say it:} a sentence of your own with \emph{der Zauberer}
  \item \emph{Recall:} say the German for the fable, then the German for the witch, then say \emph{der Zauberer} again
\end{itemize}

\begin{tcolorbox}[breakable,colback=teal!4,colframe=teal!35!black,title={Before you move on}]
What does der Zauberer mean? (\textquotedblleft{}The wizard, the magician\textquotedblright{}.) Say it once more.
\end{tcolorbox}
